\documentclass[11pt]{article}

% Page layout and typography
\usepackage[margin=1in]{geometry}
\usepackage{lmodern}
\usepackage[T1]{fontenc}
\usepackage[utf8]{inputenc}

% Figures, tables, math, and utilities
\usepackage{graphicx}
\usepackage{booktabs}
\usepackage{amsmath}
\usepackage{siunitx}
\usepackage{float}
\usepackage{enumitem}

% Hyperlinks and colors
\usepackage{xcolor}
\usepackage{hyperref}
\hypersetup{
  colorlinks=true,
  linkcolor=blue,
  urlcolor=blue,
  citecolor=blue
}

% Helper to include figures only if they exist (avoids LaTeX errors when images are not yet generated)
\newcommand{\includepng}[2][0.9\textwidth]{%
  \IfFileExists{#2}{\includegraphics[width=#1]{#2}}{\fbox{\parbox{0.8\textwidth}{Missing figure (run the notebook to generate images in \texttt{output/}).}}}%
}

\title{Pokémon Base Stat Analysis and Visual Storytelling}
\author{Wes Griffin \and Isa Wesolowski}
\date{Period 3 \quad \today}

\begin{document}

\maketitle

\begin{abstract}
We analyze base stats for \(1{,}025\) Pokémon from Kaggle, quantifying differences across generations and between Legendary/Mythical and Non-Legendary cohorts. We observe substantially higher total stats for the Legendary/Mythical cohort and modest evidence of generation-specific shifts. We close with limitations and directions for future work.
\end{abstract}

\section*{Introduction}
Pokémon designs have evolved across nine generations, blending archetypes (sweepers, walls, generalists) and type identities. Our goal is to tell a concise, visual story about how base stats vary across eras and types, how rarity (Legendary/Mythical) relates to power, and what simple models can learn from these attributes. We move through the results as a sequence of questions: who is in the dataset, how stats relate, where trade-offs appear, whether there is evidence of power creep, and which types stand out. Each plot is introduced with a short motivation and followed by a key takeaway.

\section*{Research Questions}
We study two complementary Kaggle datasets to answer the following:
\begin{enumerate}[leftmargin=*]
  \item \textbf{Dataset A (1,025 Pokémon).} How do base stats differ across generations and between Legendary/Mythical and Non-Legendary Pokémon, and what patterns emerge across primary/secondary types? We quantify correlations, cohort differences, and potential power creep across eras, and design informative visualizations (\autoref{fig:counts}, \autoref{fig:corr}, \autoref{fig:legendary}, \autoref{fig:power-mean}, \autoref{fig:power-violin}).
  \item \textbf{Dataset B (Rounak Banik).} Inspired by the dataset author's prompts, we investigate:
  \begin{enumerate}[label=\alph*)]
    \item Is it possible to build a classifier to identify legendary Pokémon? (\autoref{fig:legendary-logit}, \autoref{fig:legendary-rf})
    \item How do height and weight correlate with base stats? (\autoref{fig:hw-corr})
    \item What factors relate to Experience Growth and Egg Steps? Are they correlated with other attributes?
    \item Which primary type is strongest overall, and which is weakest? (\autoref{fig:type-strength}, \autoref{fig:type-def}, \autoref{fig:type-composite})
    \item Which primary type is most likely to be legendary? (\autoref{fig:legendary-by-type})
    \item Can we construct a ``dream team'' of six Pokémon that maximizes offensive power while remaining resilient? (\autoref{fig:dream-team})
  \end{enumerate}
\end{enumerate}

\section*{Methods}
\subsection*{Datasets}
Primary dataset: \emph{Data of 1,025 Pokémons} (Kaggle).\\
\textbf{Source:} hamzacyberpatcher / Data of 1010 Pokémons\\
\textbf{URL:} \url{https://www.kaggle.com/datasets/hamzacyberpatcher/data-of-1010-pokemons}\\
\textbf{Access date:} \today\\
\textbf{License:} See dataset page for licensing terms.

Secondary dataset: \emph{Pokemon} (Kaggle).\\
\textbf{Source:} rounakbanik / Pokemon\\
\textbf{URL:} \url{https://www.kaggle.com/datasets/rounakbanik/pokemon}\\
\textbf{Access date:} \today\\
\textbf{License:} See dataset page for licensing terms.

\paragraph{Key columns} id, name, rank, generation, evolves\_from, type1, type2, hp, atk, def, spatk, spdef, speed, total, height, weight, abilities, desc.

\paragraph{Limitations} type2 is missing for many entries; descriptions for some Pokémon are unavailable or abbreviated. No mega evolutions are included per dataset notes.

\subsection*{Data Hygiene}
\textbf{Dataset A.} We standardize fields and construct analysis helpers: (i) map generation strings to numeric (\texttt{generation\_num}); (ii) normalize \texttt{rank} to lowercase and flag \texttt{is\_legendary} (Legendary or Mythical) and \texttt{is\_baby}; (iii) preserve \texttt{type2} and add \texttt{type2\_filled}=``none'' where missing; (iv) audit missingness and duplicates (none in \texttt{id}, \texttt{type2} is the only sparse field). We also cast common categoricals and keep original text columns for later reporting.

\textbf{Dataset B.} We align naming conventions (e.g., \texttt{attack}\,$\to$\,\texttt{atk}, \texttt{sp\_attack}\,$\to$\,\texttt{spatk}, \texttt{pokedex\_number}\,$\to$\,\texttt{id}), compute \texttt{total} if absent, and ensure booleans for \texttt{is\_legendary}. When provided, we keep metric units (\texttt{height\_m}, \texttt{weight\_kg}) and the type-effectiveness columns \texttt{against\_*}.

\subsection*{Feature Engineering}
Guided by competitive archetypes and interpretability, we define:
\begin{itemize}[leftmargin=*]
  \item \textbf{Offense Index} \((\texttt{atk}+\texttt{spatk})/2\) and \textbf{Defense Index} \((\texttt{def}+\texttt{spdef})/2\) to summarize attacking and defensive prowess.
  \item \textbf{Speed-to-Power Ratio} \(\texttt{speed}/\texttt{total}\) to capture glass-cannon vs bulky tendencies.
  \item \textbf{Era Index} (\texttt{era\_total\_z}): z-score of \texttt{total} within each generation for cross-era comparability.
  \item \textbf{BMI-like} using \texttt{weight\_kg}/\texttt{height\_m}\(^2\) (Dataset B) as a rough physicality indicator.
  \item \textbf{Bulk Index} \(\sqrt{\texttt{hp}\cdot\texttt{def}\cdot\texttt{spdef}}\) to reflect survivability.
  \item \textbf{Offense Bias} \(\frac{\texttt{offense}-\texttt{defense}}{\texttt{offense}+\texttt{defense}}\) to position Pokémon on an offense--defense spectrum.
  \item \textbf{Role Tags} (heuristics): \emph{sweeper} (high speed and offense), \emph{wall} (top-quintile bulk and solid defenses), \emph{attacker} (offense-biased), \emph{tank} (defense-biased), else \emph{balanced}. See \autoref{fig:roles}.
\end{itemize}

\subsection*{Modeling and Evaluation (Dataset B)}
For legendary classification we use scikit-learn pipelines with a \texttt{ColumnTransformer} (numeric passthrough, one-hot for types/generation), splitting data with stratified 75/25 train/test. We train (i) logistic regression (max\_iter=1000) and (ii) a random forest (n=400, fixed seed). We report precision/recall/F1 and confusion matrices (\autoref{fig:legendary-logit}, \autoref{fig:legendary-rf}).

For type strength, we compute (i) mean \texttt{total} by primary type and (ii) average defensive vulnerability using \texttt{against\_*} multipliers; we combine their ranks into a composite score (\autoref{fig:type-strength}, \autoref{fig:type-def}, \autoref{fig:type-composite}). Height/weight relationships and egg-step analyses use correlation heatmaps and regression plots (\autoref{fig:hw-corr}).

For the dream team, we apply a greedy heuristic: sort by a composite score \texttt{total + 0.5*offense\_index + 0.3*speed} while enforcing diverse primary types; visualize profiles with radar small multiples (\autoref{fig:dream-team}).

\section*{Results and Discussion}
We proceed as a sequence of short questions, each paired with a figure. This interleaved format emphasizes concrete takeaways over exhaustive tables.

\paragraph{Who is in the dataset?} We start with composition by generation and type to anchor later comparisons.
\begin{figure}[H]
  \centering
  \includepng{output/generation_counts.png}
  \caption{Pokémon counts by generation.}
  \label{fig:counts}
\end{figure}
\begin{figure}[H]
  \centering
  \includepng{output/type_by_generation_heatmap.png}
  \caption{Primary type by generation (counts).}
  \label{fig:type-by-generation}
\end{figure}
\textbf{Takeaway.} The dataset spans all eras with varying type mixes across generations, reflecting design shifts (e.g., more Steel and Dark after their introduction).

\paragraph{How do base stats relate?} Correlations summarize how components contribute to overall power.
\begin{figure}[H]
  \centering
  \includepng{output/correlation_heatmap.png}
  \caption{Correlation heatmap of base stats.}
  \label{fig:corr}
\end{figure}
\textbf{Takeaway.} \texttt{total} tracks all components strongly; offense and defense pairs are moderately correlated, while \texttt{speed} is more distinct.

\paragraph{Are Legendary/Mythical Pokémon stronger?} We compare distributions of \texttt{total}.
\begin{figure}[H]
  \centering
  \includepng{output/legendary_box_total.png}
  \caption{Total stat distribution: Legendary/Mythical vs Non-Legendary.}
  \label{fig:legendary}
\end{figure}
\textbf{Takeaway.} Legendary/Mythical Pokémon concentrate at substantially higher \texttt{total} with heavier upper tails.

\paragraph{Where are trade-offs visible?} Attack vs Speed illustrates archetypes and outliers.
\begin{figure}[H]
  \centering
  \includepng{output/atk_speed_scatter.png}
  \caption{Attack vs Speed by primary type (size = total) with annotations.}
  \label{fig:scatter}
\end{figure}
\textbf{Takeaway.} A frontier of fast glass-cannons and slower heavy-hitters emerges, with type-specific clusters.

\paragraph{Has power increased over time?} We inspect means with confidence intervals and full distributions by generation.
\begin{figure}[H]
  \centering
  \includepng{output/power_creep_mean_ci.png}
  \caption{Mean total by generation with 95\% confidence intervals.}
  \label{fig:power-mean}
\end{figure}
\begin{figure}[H]
  \centering
  \includepng{output/power_creep_violin.png}
  \caption{Distribution of total by generation (violin with quartiles).}
  \label{fig:power-violin}
\end{figure}
\textbf{Takeaway.} Power creep is modest but present: later generations show slightly higher means and heavier upper tails.

\paragraph{What archetypes do types embody?} Averaging offense and defense indices by primary type reveals strategic identities.
\begin{figure}[H]
  \centering
  \includepng{output/type_centroids_off_vs_def.png}
  \caption{Primary-type centroids in offense vs defense space.}
  \label{fig:type-centroids}
\end{figure}
\textbf{Takeaway.} \emph{Dragon} and \emph{Steel} skew toward high composite power; \emph{Bug} and \emph{Normal} trend lower.

\paragraph{Can two dimensions summarize most variation?} A PCA view contextualizes overall power and offense/defense balance.
\begin{figure}[H]
  \centering
  \includepng{output/pca_base_stats.png}
  \caption{PCA of base stats colored by primary type and shaped by legendary status.}
  \label{fig:pca}
\end{figure}
\textbf{Takeaway.} The first axis tracks overall power; the second separates offense- and defense-biased designs, with legendaries occupying the high-power region.

\paragraph{Which roles are common?} Heuristic role tags illustrate population mix across tactical archetypes.
\begin{figure}[H]
  \centering
  \includepng[0.7\textwidth]{output/role_tag_counts.png}
  \caption{Counts of derived role tags.}
  \label{fig:roles}
\end{figure}
\textbf{Takeaway.} Balanced designs dominate; extreme sweepers and walls are comparatively rare.

\paragraph{How do height and weight relate to stats?} On the secondary dataset, physical attributes correlate differently with components.
\begin{figure}[H]
  \centering
  \includepng{output/height_weight_corr_heatmap.png}
  \caption{Correlation heatmap: height/weight and base stats.}
  \label{fig:hw-corr}
\end{figure}
\textbf{Takeaway.} Height/weight correlate positively with bulk-related stats and are near-zero for \texttt{speed}.

\paragraph{Which types are strongest and most resilient?} We examine mean power, defensive vulnerability, and their composite.
\begin{figure}[H]
  \centering
  \includepng{output/type_strength_mean_total.png}
  \caption{Mean total by primary type (secondary dataset).}
  \label{fig:type-strength}
\end{figure}
\begin{figure}[H]
  \centering
  \includepng{output/type_defensive_vulnerability.png}
  \caption{Average defensive vulnerability by primary type (lower is better).}
  \label{fig:type-def}
\end{figure}
\begin{figure}[H]
  \centering
  \includepng{output/type_composite_strength.png}
  \caption{Composite type strength combining power and defensive vulnerability.}
  \label{fig:type-composite}
\end{figure}
\textbf{Takeaway.} Types like \emph{Dragon}, \emph{Steel}, and \emph{Psychic} rank highly on power; defensive resilience elevates \emph{Steel} in the composite.

\paragraph{Which types are more likely to be legendary?}
\begin{figure}[H]
  \centering
  \includepng[0.8\textwidth]{output/legendary_share_by_type.png}
  \caption{Legendary share by primary type.}
  \label{fig:legendary-by-type}
\end{figure}
\textbf{Takeaway.} \emph{Dragon} and \emph{Psychic} have higher legendary prevalence; \emph{Bug} and \emph{Normal} lower, with small counts in some categories.

\paragraph{Can we identify legendaries from stats?} Simple models capture much of the signal.
\begin{figure}[H]
  \centering
  \includepng[0.6\textwidth]{output/legendary_clf_logit_cm.png}
  \caption{Legendary classifier (logistic regression) confusion matrix.}
  \label{fig:legendary-logit}
\end{figure}
\begin{figure}[H]
  \centering
  \includepng[0.6\textwidth]{output/legendary_clf_rf_cm.png}
  \caption{Legendary classifier (random forest) confusion matrix.}
  \label{fig:legendary-rf}
\end{figure}
\textbf{Takeaway.} Logistic regression and random forest separate classes well; \texttt{total} and its components dominate feature importance.

\paragraph{Can we build a strong, diverse team?} A greedy heuristic balances power, offense, speed, and type diversity.
\begin{figure}[H]
  \centering
  \includepng{output/dream_team_radars.png}
  \caption{Dream team stat profiles (heuristic selection).}
  \label{fig:dream-team}
\end{figure}
\textbf{Takeaway.} Enforcing type diversity while scoring power and speed yields a robust lineup with complementary profiles.

\section*{Limitations and Future Work}
Key limitations include missing secondary types for roughly half the entries in Dataset A and potential dataset omissions (e.g., mega evolutions). Dataset B blends stat and effectiveness data but inherits scraping noise and categorical conventions (e.g., experience growth groups). Future work: (i) incorporate complete type-effectiveness matrices and movesets; (ii) model power creep with hierarchical time-aware models; (iii) expand role-tagging with clustering; (iv) optimize dream-team selection with integer programming under coverage and synergy constraints.

\section*{Reproducibility}
Code is provided in the notebook (\texttt{main.ipynb}). Figures and tables are saved programmatically to \texttt{output/}. Environment: Python \\texttt{>= 3.12} with \texttt{pandas}, \texttt{seaborn}, \texttt{matplotlib}, \texttt{scikit-learn}, and \texttt{scipy}. To set up the environment, run \texttt{uv sync}. To regenerate figures, execute the notebook top-to-bottom; this will populate \texttt{output/} and enable LaTeX to include the images referenced in this paper.

\section*{References}
\begin{enumerate}[leftmargin=*]
  \item hamzacyberpatcher. \emph{Data of 1010 Pokémons}. Kaggle. Available at: \url{https://www.kaggle.com/datasets/hamzacyberpatcher/data-of-1010-pokemons}. Accessed: \today.
  \item Rounak Banik. \emph{Pokemon}. Kaggle. Available at: \url{https://www.kaggle.com/datasets/rounakbanik/pokemon}. Accessed: \today.
\end{enumerate}

\end{document}


